% archon:covers MR2223407ConstructionHilbertQuot/QuotConstruction.lean
\chapter{Construction of Quot Schemes}
\label{chap:quot-construction}

This chapter carries out the construction of \S5 of Nitsure: the
representability of the projective Quot functor, in the form proved by
Altman and Kleiman. The goal theorem is
\cref{thm:altman-kleiman-quot}; everything in the chapters on functors,
regularity, semicontinuity and base change, and flattening stratification
feeds into its proof.

\bigskip
\section*{Notions of Projectivity}
\medskip

\begin{definition}\leanok
[Projective morphism]
  \label{def:projective-morphism}
  \lean{MR2223407ConstructionHilbertQuot.ProjectiveMorphism}
  \uses{def:projective-space, def:very-ample, def:relative-proj}
  % SOURCE: Nitsure, §5, lines 2165–2173 (read from references/MR2223407-construction-hilbert-quot.tex)
  % SOURCE QUOTE: "Let $S$ be a noetherian scheme. Recall that as defined by Grothendieck,
  % a morphism $X\to S$ is called a \textbf{projective morphism}
  % if there exists a coherent sheaf $E$ on $S$, together with a
  % closed embedding of $X$ into $\PP(E) = \Proj \sym_{{\mathcal O}_S} E $
  % over $S$. Equivalently, $X\to S$ is projective when
  % it is proper and there exists a relatively very ample
  % line bundle $L$ on $X$ over $S$. These conditions are related by
  % taking $L$ to be the restriction of ${\mathcal O}_{\PP(E)}(1)$ to $X$, or
  % in the reverse direction, taking $E$ to be the direct image of $L$ on $S$."
  \textit{Source: Nitsure, §5.}
  Let $S$ be a noetherian scheme. A morphism $X \to S$ is \emph{projective}
  (in the sense of Grothendieck) if there exists a coherent sheaf $E$ on $S$
  together with a closed embedding $X \hookrightarrow \PP(E) = \Proj \Sym_{\Oo_S} E$
  over $S$. Equivalently, $X \to S$ is projective when it is proper and carries
  a relatively very ample line bundle $L$ over $S$. The two descriptions are
  matched by taking $L = \restr{\Oo_{\PP(E)}(1)}{X}$, or conversely $E = \pi_* L$.
\end{definition}

\begin{definition}\leanok
[Strongly projective morphism]
  \label{def:strongly-projective}
  \lean{MR2223407ConstructionHilbertQuot.StronglyProjective}
  \uses{def:projective-morphism, def:relative-proj}
  % SOURCE: Nitsure, §5, lines 2178–2184 (read from references/MR2223407-construction-hilbert-quot.tex)
  % SOURCE QUOTE: "A stronger version of projectivity was introduced
  % by Altman and Kleiman: a morphism
  % $X\to S$ of noetherian schemes is called \textbf{strongly projective}
  % (respectively, \textbf{strongly quasi-projective}) if
  % there exists a vector bundle $E$ on $S$ together with a closed
  % embedding (respectively, a locally closed embedding)
  % $X\subset \PP(E)$ over $S$."
  \textit{Source: Nitsure, §5.}
  A morphism $X \to S$ of noetherian schemes is \emph{strongly projective}
  (respectively, \emph{strongly quasi-projective}) if there exists a vector
  bundle $E$ on $S$ together with a closed embedding (respectively, a locally
  closed embedding) $X \subset \PP(E)$ over $S$. This is the special case of
  \cref{def:projective-morphism} in which the coherent sheaf $E$ is required to
  be locally free.
\end{definition}

\bigskip
\section*{Main Existence Theorems}
\medskip

\begin{theorem}\leanok
[Grothendieck Quot scheme]
  \label{thm:grothendieck-quot}
  \lean{MR2223407ConstructionHilbertQuot.quot_grothendieck}
  \uses{def:quot-functor-phi, def:projective-morphism, def:points-functor}
  % SOURCE: Nitsure, §5, lines 2230–2236 (read from references/MR2223407-construction-hilbert-quot.tex)
  % SOURCE QUOTE: "{\rm (Grothendieck)}
  % Let $S$ be a noetherian scheme, $\pi: X\to S$ a projective morphism,
  % and $L$ a relatively very ample line bundle on $X$. Then for any
  % coherent ${\mathcal O}_X$-module $E$ and any polynomial $\Phi \in \Q[\lambda]$,
  % the functor $\Quot_{E/X/S}^{\Phi,L}$ is representable
  % by a projective $S$-scheme $\quot_{E/X/S}^{\Phi,L}$."
  \textit{Source: Nitsure, §5, Theorem (Grothendieck).}
  Let $S$ be a noetherian scheme, $\pi : X \to S$ a projective morphism, and
  $L$ a relatively very ample line bundle on $X$. Then for any coherent
  $\Oo_X$-module $E$ and any polynomial $\Phi \in \QQ[\lambda]$, the functor
  $\Quotf_{E/X/S}^{\Phi,L}$ is representable by a projective $S$-scheme
  $\quotsch_{E/X/S}^{\Phi,L}$.
\end{theorem}

% SOURCE QUOTE PROOF: "We now describe how to get Theorem \ref{Grothendieck Quot}
% from the above proof. As $S$ is noetherian, we can find a common
% $m$ such that given any field-valued point
% $s: \spec k\to S$ and a coherent quotient $q: E_s \to \F$ on $X_s$
% with Hilbert polynomial $\Phi$, the sheaves $E_s(r)$, $\F(r)$, $\G(r)$
% (where $\G = \ker(q)$) are generated by global sections and all their higher
% cohomologies vanish, whenever $r\ge m$. ... Thus, we get
% $\quot_{E/X/S}^{\Phi,L}$ as a projective scheme over $S$."
\begin{proof}
  \uses{thm:altman-kleiman-quot, thm:mumford-regularity, thm:grassmannian-represents, thm:flattening-stratification, thm:valuative-criterion-properness}
  This is the noetherian form of the Altman--Kleiman theorem with the strong
  projectivity hypothesis weakened to projectivity. Since $S$ is noetherian one
  may, as in \cref{thm:altman-kleiman-quot}, choose a single integer $m$ (by
  Castelnuovo--Mumford regularity, \cref{thm:mumford-regularity}, and
  semicontinuity) so that for every field-valued point $s : \spec k \to S$ and
  every coherent quotient $q : E_s \to \Ff$ of Hilbert polynomial $\Phi$ the
  sheaves $E_s(r)$, $\Ff(r)$ and $\Gg(r) = \ker(q)(r)$ are globally generated
  with vanishing higher cohomology for $r \ge m$. As in
  \cref{lem:quot-embed-grassmannian} this gives an injection of
  $\Quotf_{E/X/S}^{\Phi,L}$ into the Grassmannian functor of $\pi_* E(r)$. Here
  $\pi_* E(r)$ is coherent but need not be a quotient of a vector bundle, so the
  Grassmannian $\grasssch(\pi_* E(r), \Phi(r))$ is projective over $S$ but not
  necessarily strongly projective. Flattening stratification
  (\cref{thm:flattening-stratification}), performed over an affine open cover of
  $S$, cuts out a locally closed subscheme representing
  $\Quotf_{E/X/S}^{\Phi,L}$, which is closed by the valuative criterion
  (\cref{thm:valuative-criterion-properness}). Hence $\quotsch_{E/X/S}^{\Phi,L}$
  is projective over $S$.
\end{proof}

\begin{theorem}\leanok
[Altman--Kleiman Quot scheme]
  \label{thm:altman-kleiman-quot}
  \lean{MR2223407ConstructionHilbertQuot.quot_altmanKleiman}
  \uses{def:quot-functor-phi, def:strongly-projective, lem:quot-reduction, thm:mumford-regularity, thm:flattening-stratification, thm:grassmannian-represents, lem:quot-embed-grassmannian, lem:quot-properness, def:points-functor}
  % SOURCE: Nitsure, §5, lines 2249–2261 (read from references/MR2223407-construction-hilbert-quot.tex)
  % SOURCE QUOTE: "{\rm (Altman-Kleiman)}
  % Let $S$ be a noetherian scheme, $X$ a closed subscheme of
  % $\PP(V)$ for some vector bundle $V$ on $S$, $L = {\mathcal O}_{\PP(V)}(1)|_X$,
  % $E$ a coherent quotient sheaf of $\pi^*(W)(\nu)$ where
  % $W$ is a vector bundle on $S$ and $\nu$ is an integer, and
  % $\Phi \in \Q[\lambda]$.
  % Then the functor $\Quot_{E/X/S}^{\Phi,L}$ is representable
  % by a scheme $\quot_{E/X/S}^{\Phi,L}$ which can be embedded over $S$ as a
  % closed subscheme of $\PP(F)$ for some vector bundle $F$ on $S$.
  % The vector bundle $F$ can be chosen to be an exterior power of
  % the tensor product of $W$ with a symmetric powers of $V$."
  \textit{Source: Nitsure, §5, Theorem (Altman--Kleiman).}
  Let $S$ be a noetherian scheme, $X$ a closed subscheme of $\PP(V)$ for some
  vector bundle $V$ on $S$, $L = \restr{\Oo_{\PP(V)}(1)}{X}$, $E$ a coherent
  quotient sheaf of $\pi^*(W)(\nu)$ where $W$ is a vector bundle on $S$ and
  $\nu$ an integer, and $\Phi \in \QQ[\lambda]$. Then the functor
  $\Quotf_{E/X/S}^{\Phi,L}$ is representable by a scheme
  $\quotsch_{E/X/S}^{\Phi,L}$ which embeds over $S$ as a closed subscheme of
  $\PP(F)$ for some vector bundle $F$ on $S$. The bundle $F$ may be taken to be
  an exterior power of the tensor product of $W$ with a symmetric power of $V$,
  namely $F = \bigwedge^{\Phi(r)}(W \ot_{\Oo_S} \Sym^r V)$ for a suitable $r$.
\end{theorem}

% SOURCE QUOTE PROOF: "The rest of this section is devoted to proving
% Theorem \ref{Altman-Kleiman Quot}, with extra noetherian hypothesis.
% ... Hence we have shown that $\Quot^{\Phi,L}_{E/X/S}$ is represented
% by a locally closed subscheme of $\grass( W \otimes_{{\mathcal O}_S}\sym^r V, \Phi(r))$.
% ... As $S$ is noetherian and as we have already shown that
% $\quot^{\Phi,L}_{E/X/S}\to S$ is of finite type, it follows that
% $\quot^{\Phi,L}_{E/X/S}\to S$ is a proper morphism. Therefore the embedding ...
% is a closed embedding. This completes the proof of Theorem \ref{Altman-Kleiman Quot}."
\begin{proof}
  \uses{lem:quot-reduction, lem:quot-uniform-regularity, lem:quot-embed-grassmannian, lem:quot-flattening-stratum, lem:quot-properness}
  The argument is the assembly of the lemmas that follow.
  \begin{enumerate}
    \item By \cref{lem:quot-reduction} it suffices to treat the model case
      $X = \PP(V)$, $E = \pi^* W$: tensoring by $L^\nu$ reduces the general
      polynomial to a shifted one, and a surjection of coherent sheaves induces
      a closed embedding of the corresponding Quot functors, so the general
      $\quotsch_{E/X/S}^{\Phi,L}$ sits as a closed subscheme of the model one.
    \item By \cref{lem:quot-uniform-regularity}, Castelnuovo--Mumford regularity
      (\cref{thm:mumford-regularity}) supplies a single integer $m$, depending
      only on $\rank V$, $\rank W$ and $\Phi$, so that for every member of the
      family the kernel, the source and the quotient become $m$-regular; their
      direct images under $\pi$ after twisting by $r \ge m$ are then locally
      free of fixed rank with vanishing higher direct images.
    \item By \cref{lem:quot-embed-grassmannian}, sending
      $\langle \Ff, q\rangle$ to the quotient
      $\pi_* E(r) \to \pi_* \Ff(r)$ of rank $\Phi(r)$ defines an injective
      morphism of functors
      $\alpha : \Quotf_{E/X/S}^{\Phi,L} \to \grassf(W \ot_{\Oo_S} \Sym^r V, \Phi(r))$,
      with $\pi_* E(r) = W \ot_{\Oo_S} \Sym^r V$; the Grassmannian functor is
      representable (\cref{thm:grassmannian-represents}).
    \item By \cref{lem:quot-flattening-stratum}, the image of $\alpha$ is the
      flattening stratum (\cref{thm:flattening-stratification}) for the
      universal quotient over the Grassmannian for Hilbert polynomial $\Phi$,
      giving a locally closed subscheme that represents
      $\Quotf_{E/X/S}^{\Phi,L}$ and thus a locally closed embedding
      $\quotsch_{E/X/S}^{\Phi,L} \subset \PP(\bigwedge^{\Phi(r)}(W \ot_{\Oo_S}\Sym^r V))$.
      In particular $\quotsch_{E/X/S}^{\Phi,L} \to S$ is separated and of finite
      type.
    \item By \cref{lem:quot-properness}, the valuative criterion shows
      $\quotsch_{E/X/S}^{\Phi,L} \to S$ is proper; being also of finite type over
      the noetherian $S$, the locally closed embedding above is in fact a closed
      embedding into $\PP(F)$. This completes the proof.\qedhere
  \end{enumerate}
\end{proof}

\begin{theorem}\leanok
[Quot scheme in projective space]
  \label{thm:quot-pn}
  \lean{MR2223407ConstructionHilbertQuot.quot_projectiveSpace}
  \uses{thm:altman-kleiman-quot}
  % SOURCE: Nitsure, §5, lines 2265–2273 (read from references/MR2223407-construction-hilbert-quot.tex)
  % SOURCE QUOTE: "If $S$ is a noetherian scheme, $X$ is a closed subscheme of
  % $\PP^n_S$ for some $n\ge 0$, $L = {\mathcal O}_{\PP^n_S}(1)|_X$,
  % $E$ is a coherent quotient sheaf of $\oplus^p {\mathcal O}_X(\nu)$
  % for some integers $p\ge 0$ and $\nu$, and $\Phi \in \Q[\lambda]$,
  % then the the functor $\Quot_{E/X/S}^{\Phi,L}$ is representable
  % by a scheme $\quot_{E/X/S}^{\Phi,L}$ which can be embedded over $S$ as a
  % closed subscheme of $\PP^r_S$ for some $r\ge 0$."
  \textit{Source: Nitsure, §5, Theorem (trivial-bundle case).}
  Let $S$ be a noetherian scheme, $X$ a closed subscheme of $\PP^n_S$ for some
  $n \ge 0$, $L = \restr{\Oo_{\PP^n_S}(1)}{X}$, $E$ a coherent quotient sheaf of
  $\oplus^p \Oo_X(\nu)$ for integers $p \ge 0$ and $\nu$, and
  $\Phi \in \QQ[\lambda]$. Then $\Quotf_{E/X/S}^{\Phi,L}$ is representable by a
  scheme $\quotsch_{E/X/S}^{\Phi,L}$ which embeds over $S$ as a closed subscheme
  of $\PP^r_S$ for some $r \ge 0$.
\end{theorem}

\begin{proof}
  \uses{thm:altman-kleiman-quot}
  Apply \cref{thm:altman-kleiman-quot} with $V = \Oo_S^{\,n+1}$ and
  $W = \Oo_S^{\,p}$ both trivial. Then $\PP(V) = \PP^n_S$, the sheaf
  $\pi^*(W)(\nu) = \oplus^p \Oo_X(\nu)$, and the closed embedding target
  $\PP(\bigwedge^{\Phi(r)}(W \ot_{\Oo_S}\Sym^r V))$ becomes a trivial projective
  bundle $\PP^r_S$ with $r + 1 = \rank \bigwedge^{\Phi(r)}(W \ot_{\Oo_S}\Sym^r V)$.
\end{proof}

\bigskip
\section*{Reduction to the case of $\Quotf_{\pi^*W/\PP(V)/S}^{\Phi,L}$}
\medskip

\begin{lemma}\leanok
[Reduction to the model case]
  \label{lem:quot-reduction}
  \lean{MR2223407ConstructionHilbertQuot.quot_reduction}
  \uses{def:quot-functor-phi}
  % SOURCE: Nitsure, §5, lines 2299–2310 (read from references/MR2223407-construction-hilbert-quot.tex)
  % SOURCE QUOTE: "(i) Let $\nu$ be any integer.
  % Then tensoring by $L^{\nu}$ gives an isomorphism of
  % functors from $\Quot_{E/X/S}^{\Phi,L}$ to
  % $\Quot_{E(\nu)/X/S}^{\Psi,L}$ where the polynomial
  % $\Psi \in \Q[\lambda]$ is defined by $\Psi(\lambda) = \Phi(\lambda + \nu)$.
  % (ii) Let $\phi : E \to G$ be a surjective homomorphism of
  % coherent sheaves on $X$. Then the corresponding natural transformation
  % $\Quot_{G/X/S}^{\Phi,L} \to \Quot_{E/X/S}^{\Phi,L}$
  % is a closed embedding."
  \textit{Source: Nitsure, §5, Lemma (reduction).}
  Let $\pi : X \to S$ be projective with relatively very ample $L$.
  \begin{enumerate}
    \item For any integer $\nu$, tensoring by $L^\nu$ gives an isomorphism of
      functors $\Quotf_{E/X/S}^{\Phi,L} \cong \Quotf_{E(\nu)/X/S}^{\Psi,L}$,
      where $\Psi \in \QQ[\lambda]$ is defined by $\Psi(\lambda) = \Phi(\lambda + \nu)$.
    \item For any surjective homomorphism $\phi : E \to G$ of coherent sheaves
      on $X$, the induced natural transformation
      $\Quotf_{G/X/S}^{\Phi,L} \to \Quotf_{E/X/S}^{\Phi,L}$ is a closed embedding.
  \end{enumerate}
  Consequently, to prove \cref{thm:altman-kleiman-quot} it suffices to treat the
  model case $X = \PP(V)$ and $E = \pi^* W$ for vector bundles $V, W$ on $S$:
  every $\quotsch_{E/X/S}^{\Phi,L}$ then arises as a closed subscheme of
  $\quotsch_{\pi^*W/\PP(V)/S}^{\Phi,L}$.
\end{lemma}

% SOURCE QUOTE PROOF: "The statement (i) is obvious. The statement (ii) just says that
% given any locally noetherian scheme $T$ and a family
% $\langle \F,q\rangle \in \Quot_{E/X/S}^{\Phi,L}(T)$, there exists
% a closed subscheme $T' \subset T$ with the following universal property:
% for any locally noetherian scheme $U$ and a morphism $f: U\to T$,
% the pulled back homomorphism of ${\mathcal O}_{X_U}$-modules $q_U : E_U \to \F_U$
% factors via the pulled back homomorphism $\phi_U : E_U \to G_U$
% if and only if $U \to T$ factors via $T'\hra T$.
% This is satisfied by taking $T'$ to be the vanishing scheme
% for the composite homomorphism $\ker(\phi) \hra  E \stackrel{q}{\to} \F$
% of coherent sheaves on $X_T$ ..., which makes sense here as both $\ker(\phi)$
% and $\F$ are coherent on $X_T$ and $\F$ is flat over $T$."
\begin{proof}
  \uses{def:quot-functor-phi}
  Part (i) is immediate: $L^\nu$-twisting is an equivalence on coherent
  quotients, and a fibrewise twist shifts the Hilbert polynomial of a quotient
  by $\nu$, so $\langle \Ff, q\rangle \mapsto \langle \Ff(\nu), q(\nu)\rangle$
  is the claimed isomorphism with $\Psi(\lambda) = \Phi(\lambda + \nu)$.

  Part (ii) is the statement that for a locally noetherian $T$ and a family
  $\langle \Ff, q\rangle \in \Quotf_{E/X/S}^{\Phi,L}(T)$ there is a closed
  subscheme $T' \subset T$ such that a morphism $U \to T$ factors through $T'$
  if and only if the pulled-back quotient $q_U : E_U \to \Ff_U$ factors through
  $\phi_U : E_U \to G_U$. Take $T'$ to be the vanishing scheme of the composite
  $\ker(\phi) \hookrightarrow E \xrightarrow{q} \Ff$ of coherent sheaves on
  $X_T$; this is well-defined because $\ker(\phi)$ and $\Ff$ are coherent on
  $X_T$ and $\Ff$ is flat over $T$, and it manifestly has the required universal
  property. Hence the natural transformation is a closed embedding.
\end{proof}

\bigskip
\section*{Use of $m$-Regularity}
\medskip

\begin{lemma}\leanok
[Uniform $m$-regularity of the family]
  \label{lem:quot-uniform-regularity}
  \lean{MR2223407ConstructionHilbertQuot.quot_uniformRegularity}
  \uses{thm:mumford-regularity, lem:castelnuovo, def:hilbert-polynomial, thm:base-change-flat}
  % SOURCE: Nitsure, §5, lines 2358–2387 (read from references/MR2223407-construction-hilbert-quot.tex)
  % SOURCE QUOTE: "It follows from Theorem \ref{Mumford}
  % that given any $\Phi \in \Q[\lambda]$, there exists an integer
  % $m$ which depends only on $\rank(V)$, $\rank(W)$ and $\Phi$,
  % such that for any field $k$ and a $k$-valued point $s$ of $S$,
  % the sheaf $E_s$ on $\PP(V)_s$ is $m$-regular, and for any
  % coherent quotient $q: E_s \to \F$ on $\PP(V)_s$ with
  % Hilbert polynomial $\Phi$, the sheaf $\F$ and the kernel sheaf
  % $\G\subset E_s$ of $q$ are both $m$-regular. ...
  % The sheaves ${\pi_T}_*\G(r)$, ${\pi_T}_*E_T(r)$,
  % ${\pi_T}_*\F(r)$ are locally free of fixed ranks
  % determined by the data $n$, $p$, $r$, and $\Phi$, ...
  % and the higher direct images
  % $R^i{\pi_T}_*\G(r)$, $R^i{\pi_T}_*E_T(r)$,
  % $R^i{\pi_T}_*\F(r)$ are zero, for all $r\ge m$ and $i\ge 1$."
  \textit{Source: Nitsure, §5, ``Use of $m$-Regularity''.}
  Work in the model case $E = \pi^* W$ on $X = \PP(V)$ with
  $L = \Oo_{\PP(V)}(1)$, set $n = \rank V - 1$ and $p = \rank W$. By
  \cref{thm:mumford-regularity} there is a single integer $m$, depending only on
  $\rank V$, $\rank W$ and $\Phi$, such that for every field-valued point $s$ of
  $S$ the sheaf $E_s$ is $m$-regular and, for every coherent quotient
  $q : E_s \to \Ff$ on $\PP(V)_s$ of Hilbert polynomial $\Phi$, both $\Ff$ and
  the kernel $\Gg = \ker(q)$ are $m$-regular. Consequently, for any $S$-scheme
  $T$ and any $T$-flat coherent quotient $q : E_T \to \Ff$ of Hilbert
  polynomial $\Phi$ with kernel $\Gg$, and for all $r \ge m$: the sheaves
  $\pi_{T*}\Gg(r)$, $\pi_{T*}E_T(r)$, $\pi_{T*}\Ff(r)$ are locally free of fixed
  ranks determined by $n, p, r, \Phi$; the natural maps
  $\pi_T^*\pi_{T*}(-) \to (-)$ are surjective; and the higher direct images
  $R^i\pi_{T*}(-)$ vanish for $i \ge 1$. In particular $\pi_{T*}\Ff(r)$ has rank
  $\Phi(r)$ and $\pi_*E(r) = W \ot_{\Oo_S} \Sym^r V$.
\end{lemma}

% SOURCE QUOTE PROOF: "In particular, it follows from the Castelnuovo Lemma \ref{Castelnuovo}
% that for $r\ge m$, all cohomologies
% $H^i(X_s,E_s(r))$, $H^i(X_s,\F(r))$, and $H^i(X_s,\G(r))$
% are zero for $i\ge 1$, and $H^0(X_s,E_s(r))$, $H^0(X_s,\F(r))$, and
% $H^0(X_s,\G(r))$ are generated by their global sections.
% From the above it follows by Theorem \ref{Base-change for flat sheaves} that
% if $T$ is an $S$-scheme and $q :E_T \to \F$ is a $T$-flat coherent quotient
% with Hilbert polynomial $\Phi$, then ... [(*) and (**) hold]."
\begin{proof}
  \uses{thm:mumford-regularity, lem:castelnuovo, thm:cohomology-base-change, def:hilbert-polynomial}
  Fibrewise, $\PP(V)_s \cong \PP^n_k$ and $E_s \cong \oplus^p \Oo_{\PP(V)_s}$,
  so Mumford's bound (\cref{thm:mumford-regularity}) applied to the kernels of
  quotients with the fixed Hilbert polynomial $\Phi$ yields one integer $m$
  uniform over all such fibres. By the Castelnuovo Lemma
  (\cref{lem:castelnuovo}), $m$-regularity gives, for $r \ge m$ and $i \ge 1$,
  the vanishing $H^i(X_s, -(r)) = 0$ and the global generation of
  $H^0(X_s, -(r))$ for each of $E_s$, $\Ff$, $\Gg$. Cohomology and base change
  for flat sheaves (\cref{thm:cohomology-base-change}) then propagates this over
  any base $T$: the direct images $\pi_{T*}(-)(r)$ are locally free of the rank
  computed fibrewise from $n, p, r, \Phi$, the higher direct images vanish, and
  the canonical evaluation maps $\pi_T^*\pi_{T*}(-)(r) \to (-)(r)$ are
  surjective. Twisting the short exact sequence
  $0 \to \Gg \to E_T \to \Ff \to 0$ by $r$ and applying $\pi_{T*}$ keeps it
  exact, which is the diagram used below.
\end{proof}

\bigskip
\section*{Embedding Quot into the Grassmannian}
\medskip

\begin{lemma}\leanok
[Embedding of Quot into a Grassmannian]
  \label{lem:quot-embed-grassmannian}
  \lean{MR2223407ConstructionHilbertQuot.quot_embedGrassmannian}
  \uses{lem:quot-uniform-regularity, thm:grassmannian-represents}
  % SOURCE: Nitsure, §5, lines 2413–2450 (read from references/MR2223407-construction-hilbert-quot.tex)
  % SOURCE QUOTE: "We now fix a positive integer $r$ such that $r\ge m$.
  % Note that the rank of ${\pi_T}_*\F(r)$ is $\Phi(r)$
  % and $\pi_*E(r) = W \otimes_{{\mathcal O}_S}\sym^r V$.
  % Therefore the surjective homomorphism
  % ${\pi_T}_*E_T(r) \to {\pi_T}_*\F(r)$
  % defines an element of the set
  % $\Grass( W \otimes_{{\mathcal O}_S}\sym^r V , \Phi(r))(T)$.
  % We thus get a morphism of functors
  % $\alpha : \Quot^{\Phi,L}_{E/X/S} \to
  % \Grass( W \otimes_{{\mathcal O}_S}\sym^r V, \Phi(r))$
  % It associates to $q: E_T \to \F$ the quotient
  % ${\pi_T}_*(q(r)) : {\pi_T}_*E_T(r) \to {\pi_T}_*\F(r)$.
  % The above morphism $\alpha$ is injective ..."
  \textit{Source: Nitsure, §5, ``Embedding Quot into Grassmannian''.}
  Fix an integer $r \ge m$. Since $\pi_{T*}\Ff(r)$ has rank $\Phi(r)$ and
  $\pi_*E(r) = W \ot_{\Oo_S}\Sym^r V$, the surjection
  $\pi_{T*}E_T(r) \to \pi_{T*}\Ff(r)$ defines a $T$-point of the Grassmannian
  $\grasssch(W \ot_{\Oo_S}\Sym^r V, \Phi(r))$. This gives a morphism of functors
  \[
    \alpha : \Quotf_{E/X/S}^{\Phi,L} \to \grassf(W \ot_{\Oo_S}\Sym^r V, \Phi(r)),
    \qquad \langle \Ff, q\rangle \mapsto \pi_{T*}(q(r)),
  \]
  and $\alpha$ is injective. The Grassmannian functor is representable
  (\cref{thm:grassmannian-represents}).
\end{lemma}

% SOURCE QUOTE PROOF: "The above morphism $\alpha$ is injective because the quotient
% $q: E_T \to \F$ can be recovered from
% ${\pi_T}_*(q(r)) : {\pi_T}_*E_T(r) \to {\pi_T}_*\F(r)$ as follows. ...
% Let $h$ be the composite
% ${\pi_T}^*{\pi_T}_*(\G(r))\to {\pi_T}^*{\pi_T}_*(E_T(r))\to E_T(r)$.
% As a consequence of the properties of the diagram \textbf{(**)},
% the following is a right exact sequence on $X_T$
% ${\pi_T}^*{\pi_T}_*(\G(r))\stackrel{h}{\to} E_T(r) \stackrel{q(r)}{\to} \F \to 0$
% and so $q(r) : E_T(r) \to \F(r)$ can be recovered as the cokernel of $h$.
% Finally, twisting by $-r$, we recover $q$, proving the desired injectivity ..."
\begin{proof}
  \uses{lem:quot-uniform-regularity, thm:grassmannian-represents}
  The map $\alpha$ is well-defined by \cref{lem:quot-uniform-regularity}: for
  $r \ge m$, $\pi_{T*}(q(r))$ is a locally free quotient of
  $\pi_{T*}E_T(r) = W_T \ot \Sym^r V_T$ of rank $\Phi(r)$, hence a $T$-point of
  the Grassmannian, which is representable by \cref{thm:grassmannian-represents}.
  For injectivity, the kernel inclusion $\pi_{T*}\Gg(r) \hookrightarrow
  \pi_{T*}E_T(r)$ is recovered from $\pi_{T*}(q(r))$. Let $h$ be the composite
  $\pi_T^*\pi_{T*}\Gg(r) \to \pi_T^*\pi_{T*}E_T(r) \to E_T(r)$. By the surjective
  evaluation maps and exact rows of \cref{lem:quot-uniform-regularity},
  \[
    \pi_T^*\pi_{T*}\Gg(r) \xrightarrow{\ h\ } E_T(r) \xrightarrow{\ q(r)\ } \Ff(r) \to 0
  \]
  is right exact, so $q(r)$ is the cokernel of $h$; twisting by $-r$ recovers
  $q$. Thus $\alpha$ is injective.
\end{proof}

\bigskip
\section*{Use of Flattening Stratification}
\medskip

\begin{lemma}\leanok
[The representing flattening stratum]
  \label{lem:quot-flattening-stratum}
  \lean{MR2223407ConstructionHilbertQuot.quot_flatteningStratum}
  \uses{lem:quot-embed-grassmannian, thm:flattening-stratification, lem:grassmannian-very-ample}
  % SOURCE: Nitsure, §5, lines 2460–2503 (read from references/MR2223407-construction-hilbert-quot.tex)
  % SOURCE QUOTE: "We will next prove that
  % $\alpha :\Quot^{\Phi,L}_{E/X/S} \to
  % \Grass( W \otimes_{{\mathcal O}_S}\sym^r V, \Phi(r))$
  % is relatively representable. In fact, we will show that
  % given any locally noetherian $S$-scheme $T$ and
  % a surjective homomorphism $f :  W_T \otimes_{{\mathcal O}_T}\sym^r V_T \to \J$
  % where $\J$ is a locally free ${\mathcal O}_T$-module of rank $\Phi(r)$,
  % there exists a locally closed subscheme $T'$ of
  % $T$ with the following universal property \textbf{(F)} ...
  % The existence of such a locally closed subscheme $T'$ of
  % $T$ is given by Theorem \ref{flattening stratification}, which
  % shows that $T'$ is the stratum corresponding to
  % Hilbert polynomial $\Phi$ for the
  % flattening stratification over $T$ for the sheaf $\F$ on $X_T$."
  \textit{Source: Nitsure, §5, ``Use of Flattening Stratification''.}
  The morphism $\alpha$ of \cref{lem:quot-embed-grassmannian} is relatively
  representable by locally closed immersions. Concretely, for any locally
  noetherian $S$-scheme $T$ and any surjection
  $f : W_T \ot_{\Oo_T}\Sym^r V_T \to \J$ onto a locally free $\Oo_T$-module $\J$
  of rank $\Phi(r)$, form $\Kk = \ker(f)$, the composite
  $h : \pi_Y^*\Kk_Y \to \pi_Y^*(W \ot_{\Oo_S}\Sym^r V) = \pi_Y^*\pi_{Y*}E_Y \to E_Y$
  after pullback along $\phi : Y \to T$, and the cokernel $q : E_Y \to \Ff$ of
  $h$. There is a locally closed subscheme $T' \subset T$ such that $\phi$
  factors through $T'$ if and only if $\Ff$ is $Y$-flat with fibrewise Hilbert
  polynomial $\Phi$: namely $T'$ is the stratum of Hilbert polynomial $\Phi$ in
  the flattening stratification of $\Ff$ on $X_T$
  (\cref{thm:flattening-stratification}). Taking
  $T = \grasssch(W \ot_{\Oo_S}\Sym^r V, \Phi(r))$ with its universal quotient,
  the corresponding $T'$ represents $\Quotf_{E/X/S}^{\Phi,L}$, yielding a locally
  closed embedding
  $\quotsch_{E/X/S}^{\Phi,L} \subset \PP(\bigwedge^{\Phi(r)}(W \ot_{\Oo_S}\Sym^r V))$;
  in particular $\quotsch_{E/X/S}^{\Phi,L} \to S$ is separated and of finite type.
\end{lemma}

% SOURCE QUOTE PROOF: "When we take $T$ to be $\grass( W \otimes_{{\mathcal O}_S}\sym^r V, \Phi(r))$
% with universal quotient $u : {p_G}^*E \to \U$, the corresponding
% locally closed subscheme $T'$ represents the functor $\Quot^{\Phi,L}_{E/X/S}$
% by its construction. Hence we have shown that $\Quot^{\Phi,L}_{E/X/S}$ is represented
% by a locally closed subscheme of $\grass( W \otimes_{{\mathcal O}_S}\sym^r V, \Phi(r))$.
% As $\grass( W \otimes_{{\mathcal O}_S}\sym^r V, \Phi(r))$ embeds as a
% closed subscheme of $\PP(\bigwedge ^{\Phi(r)} W \otimes_{{\mathcal O}_S}\sym^r V)$,
% we get a locally closed embedding of $S$-schemes
% $\quot^{\Phi,L}_{E/X/S} \subset \PP(\bigwedge ^{\Phi(r)} (W \otimes_{{\mathcal O}_S}\sym^r V))$.
% In particular, the morphism $\quot^{\Phi,L}_{E/X/S}\to S$ is separated and of finite type."
\begin{proof}
  \uses{lem:quot-embed-grassmannian, thm:flattening-stratification}
  Given the data $f, \J$ over $T$, the cokernel construction $q : E_Y \to \Ff$
  is functorial in $\phi : Y \to T$, and by the recovery in
  \cref{lem:quot-embed-grassmannian} a $T$-point lies in the image of $\alpha$
  exactly when this $\Ff$ is flat over $Y$ with fibrewise Hilbert polynomial
  $\Phi$. The flattening stratification theorem
  (\cref{thm:flattening-stratification}) provides, for the sheaf $\Ff$ on $X_T$,
  a locally finite stratification of $T$ by locally closed subschemes indexed by
  Hilbert polynomials, with the universal property that $\phi : Y \to T$
  pulls $\Ff$ back to a $Y$-flat sheaf of fibrewise polynomial $\Phi$ iff $\phi$
  factors through the stratum $T'$ of polynomial $\Phi$. Applying this to
  $T = \grasssch(W \ot_{\Oo_S}\Sym^r V, \Phi(r))$ and its universal quotient,
  $T'$ represents $\Quotf_{E/X/S}^{\Phi,L}$. Since the Grassmannian embeds as a
  closed subscheme of $\PP(\bigwedge^{\Phi(r)}(W \ot_{\Oo_S}\Sym^r V))$ via the
  Plücker embedding, $T' = \quotsch_{E/X/S}^{\Phi,L}$ acquires a locally closed
  embedding into that projective bundle, hence is separated and of finite type
  over $S$.
\end{proof}

\bigskip
\section*{Valuative Criterion for Properness}
\medskip

\begin{lemma}\leanok
[Properness via the valuative criterion]
  \label{lem:quot-properness}
  \lean{MR2223407ConstructionHilbertQuot.quot_properness}
  \uses{lem:quot-flattening-stratum, thm:valuative-criterion-properness}
  % SOURCE: Nitsure, §5, lines 2517–2540 (read from references/MR2223407-construction-hilbert-quot.tex)
  % SOURCE QUOTE: "The functor $\Quot_{E/X/S}^{\Phi,L}$ satisfies the following valuative
  % criterion for properness over $S$: given any discrete valuation ring
  % $R$ over $S$ with quotient field $K$, the restriction map
  % $\Quot_{E/X/S}^{\Phi,L}(\spec R) \to \Quot_{E/X/S}^{\Phi,L}(\spec K)$
  % is bijective. ...
  % As $S$ is noetherian and as we have already shown that
  % $\quot^{\Phi,L}_{E/X/S}\to S$ is of finite type,
  % it follows that $\quot^{\Phi,L}_{E/X/S}\to S$ is a proper morphism.
  % Therefore the embedding of $\quot^{\Phi,L}_{E/X/S}$ into
  % $\PP(\bigwedge ^{\Phi(r)} (W \otimes_{{\mathcal O}_S}\sym^r V))$
  % is a closed embedding."
  \textit{Source: Nitsure, §5, ``Valuative Criterion for Properness''.}
  The functor $\Quotf_{E/X/S}^{\Phi,L}$ satisfies the valuative criterion of
  properness over $S$: for every discrete valuation ring $R$ over $S$ with
  fraction field $K$, the restriction map
  \[
    \Quotf_{E/X/S}^{\Phi,L}(\spec R) \to \Quotf_{E/X/S}^{\Phi,L}(\spec K)
  \]
  is bijective. Combined with finite type over the noetherian $S$
  (\cref{lem:quot-flattening-stratum}), it follows that
  $\quotsch_{E/X/S}^{\Phi,L} \to S$ is proper, and hence the locally closed
  embedding into $\PP(\bigwedge^{\Phi(r)}(W \ot_{\Oo_S}\Sym^r V))$ is a closed
  embedding.
\end{lemma}

% SOURCE QUOTE PROOF: "Given any coherent quotient
% $q: E_K \to \F$ on $X_R$ which defines an element $\langle \F, q\rangle$
% of $\Quot_{E/X/S}^{\Phi,L}(\spec K)$.
% Let $\ov{\F}$ be the image of the composite homomorphism
% $E_R \to j_*(E_K) \to j_*\F$ where $j : X_K \hra X_R$ is the open
% inclusion. Let $\ov{q} : E_R\to \ov{\F}$ be the induced surjection.
% Then we leave it to the reader to verify that
% $\langle \F, q\rangle$ is an element of $\Quot_{E/X/S}^{\Phi,L}(\spec R)$
% which maps to $\langle \F, q\rangle$, and is the unique such element.
% (Use the basic fact that being flat over a dvr is the same as being torsion-free.)"
\begin{proof}
  \uses{lem:quot-flattening-stratum, thm:valuative-criterion-properness}
  Let $R$ be a discrete valuation ring over $S$ with fraction field $K$ and
  $j : X_K \hookrightarrow X_R$ the open inclusion. Given
  $\langle \Ff, q\rangle \in \Quotf_{E/X/S}^{\Phi,L}(\spec K)$, let
  $\overline{\Ff}$ be the image of the composite
  $E_R \to j_*(E_K) \to j_*\Ff$ and $\overline{q} : E_R \to \overline{\Ff}$ the
  induced surjection. Then $\overline{\Ff}$ is torsion-free over $R$, hence flat
  (flatness over a dvr is torsion-freeness), with the same Hilbert polynomial
  $\Phi$, so $\langle \overline{\Ff}, \overline{q}\rangle$ is an $R$-point
  restricting to $\langle \Ff, q\rangle$; uniqueness follows from the same
  torsion-free characterization. This verifies the valuative criterion, so by
  \cref{thm:valuative-criterion-properness} and finite type over the noetherian
  $S$ (\cref{lem:quot-flattening-stratum}), $\quotsch_{E/X/S}^{\Phi,L} \to S$ is
  proper. A locally closed immersion that is also proper is a closed immersion,
  giving the closed embedding into the projective bundle.
\end{proof}

\bigskip
\section*{The Version of Grothendieck}
\medskip

We close by deriving \cref{thm:grothendieck-quot} from
\cref{thm:altman-kleiman-quot}. Since $S$ is noetherian, the theory of
$m$-regularity together with semicontinuity (\cref{thm:mumford-regularity})
furnishes a common integer $m$ such that for every field-valued point
$s : \spec k \to S$ and every coherent quotient $q : E_s \to \Ff$ of Hilbert
polynomial $\Phi$, the sheaves $E_s(r)$, $\Ff(r)$ and $\Gg(r) = \ker(q)(r)$ are
globally generated with vanishing higher cohomology for $r \ge m$. As in
\cref{lem:quot-embed-grassmannian} this yields an injection of
$\Quotf_{E/X/S}^{\Phi,L}$ into the Grassmannian functor
$\grassf(\pi_* E(r), \Phi(r))$. Here $\pi_* E(r)$ is coherent but need not be a
quotient of a vector bundle, so $\grasssch(\pi_* E(r), \Phi(r))$ is projective
over $S$ but not necessarily strongly projective. Finally, flattening
stratification (\cref{thm:flattening-stratification}), carried out over an
affine open cover of $S$, gives a locally closed subscheme of
$\grasssch(\pi_* E(r), \Phi(r))$ representing $\Quotf_{E/X/S}^{\Phi,L}$, which is
a closed subscheme by the valuative criterion
(\cref{lem:quot-properness}). Thus $\quotsch_{E/X/S}^{\Phi,L}$ is a projective
$S$-scheme, proving \cref{thm:grothendieck-quot}.

\section{Auxiliary declarations}

\begin{definition}

\leanok
[Universe-lifted functor of points]
  \label{def:points-functor}
  \lean{MR2223407ConstructionHilbertQuot.pointsFunctor}
  For an object $A$ of $\Sch/S$, the \emph{functor of points} of $A$ raised one
  universe level: the representable presheaf $h_A = \Mor_S(-,A)$ post-composed
  with the universe lift, so it lands in $\mathrm{Type}\,(u+1)$. A moduli functor
  $M\colon(\Sch/S)^{\mathrm{op}}\to\mathrm{Type}\,(u+1)$ is then said to be
  \emph{represented} by $A$ exactly when $M\cong h_A$; this is the
  representability shape used in \cref{thm:grothendieck-quot} and
  \cref{thm:altman-kleiman-quot}. The level shift is forced: $\Mor$-sets live in
  $\mathrm{Type}\,u$ whereas the Quot moduli functors of this chapter are valued
  in $\mathrm{Type}\,(u+1)$, so the bare Yoneda embedding cannot be compared to
  them without the lift.
\end{definition}
